La naturaleza autorregresiva de los LLMs hace que una respuesta extensa pueda tardar varios segundos en generarse. Esperar a que Ollama termine por completo antes de enviar el bloque de texto al cliente produciría una pantalla congelada inaceptable en un entorno de tutoría interactiva. La solución es abrir una conexión de larga duración y enviar cada fragmento de texto en el instante exacto en que el modelo lo produce, de modo que el alumno ve la respuesta construirse progresivamente.

El flujo de datos del tutor socrático es estrictamente unidireccional durante la generación: el servidor empuja \textit{tokens} al cliente y el cliente no necesita enviar nada hasta que el \textit{stream} ha terminado. Esta asimetría hace que WebSockets, un protocolo \textit{full-duplex}, añada complejidad innecesaria. Requiere un \textit{protocol upgrade} explícito, mantener el estado de la conexión en el servidor y gestionar la reconexión manualmente. La alternativa de \textit{polling} periódico implicaría peticiones continuas aunque no haya datos nuevos. Server-Sent Events (SSE) se ajusta de forma natural al problema: utiliza una conexión HTTP estándar de larga duración con el tipo de contenido \texttt{text/event-stream}, el navegador la maneja de forma nativa y el servidor puede cerrarla cuando el \textit{stream} termina sin ningún protocolo adicional.

\paragraph{Backend.}

En el servidor, la respuesta no se construye en memoria para enviarse al final, sino que una función generadora asíncrona \texttt{yield}-ea fragmentos conforme los recibe de Ollama. Django envuelve ese generador en un \texttt{StreamingHttpResponse} con el tipo de contenido adecuado. Dentro del generador, el cliente asíncrono de Ollama itera sobre el \textit{stream} de \textit{chunks} y reenvía cada \textit{token} al cliente en formato SSE estándar:

\PythonCode[lst:streaming_loop]{Bucle de streaming}{Iteración \textit{token} a \textit{token} sobre el stream de Ollama con reenvío SSE inmediato (\texttt{views.py}).}{python/streaming_loop.py}{1}{9}{1}

La llamada a \texttt{asyncio.sleep(0)} tras cada \textit{yield} cede el control al \textit{event loop}, lo que permite que las tareas de moderación de la salida descritas en la \autoref{SS:DI_IMP_MOD} se ejecuten de forma concurrente sin bloquear el flujo de \textit{tokens}.

\paragraph{Frontend.}

En el cliente, la API estándar \texttt{EventSource} no es adecuada porque solo admite peticiones GET, mientras que el envío del mensaje del alumno requiere un POST con cuerpo JSON. Por ello, \texttt{chatService.js} utiliza la \textit{Fetch API} nativa, que expone el cuerpo de la respuesta como un \texttt{ReadableStream} leído de forma incremental. Un \texttt{TextDecoder} con el flag \texttt{stream: true} convierte los bytes en texto UTF-8 sin descartar secuencias incompletas entre \textit{chunks}. Cada línea con el prefijo \texttt{data:} se parsea como JSON y se despacha al callback correspondiente según el tipo de evento.

Antes de iniciar el \textit{stream}, la \textit{action} \texttt{sendMessage} del store de Pinia inserta un objeto \textit{placeholder} en el array de mensajes con \texttt{content} vacío y un flag \texttt{\_isStreaming: true}. A partir de ese momento, cada \textit{token} entrante se acumula en \texttt{streamBuffer} y se asigna directamente al campo \texttt{content} del \textit{placeholder}. Dado que \texttt{messages} es un array reactivo de Pinia, cada asignación dispara automáticamente el ciclo de re-renderizado de Vue 3 y el resultado visual es el efecto de máquina de escribir que percibe el alumno.
